\documentclass[11pt]{article}
\usepackage[T1]{fontenc}
\usepackage{lmodern,amsmath,amssymb,amsthm,mathtools,mathrsfs}
\usepackage[nopatch=footnote]{microtype}
\usepackage[margin=1in]{geometry}
\usepackage[colorlinks=true,linkcolor=black,citecolor=black,urlcolor=black]{hyperref}
\usepackage{enumitem,needspace}
\hypersetup{
 pdftitle={A proof of Nivat's conjecture},
 pdfauthor={Pranava Kumar},
 pdfsubject={Rectangular pattern complexity and periodicity}
}
\numberwithin{equation}{section}
\setlength{\emergencystretch}{1em}
\clubpenalty=10000
\widowpenalty=10000
\displaywidowpenalty=10000
\frenchspacing
\setlist{itemsep=2pt,topsep=5pt}
\newtheorem{theorem}{Theorem}[section]
\newtheorem{lemma}[theorem]{Lemma}
\newtheorem{proposition}[theorem]{Proposition}
\newtheorem{corollary}[theorem]{Corollary}
\theoremstyle{definition}
\newtheorem{definition}[theorem]{Definition}
\theoremstyle{remark}
\newtheorem{remark}[theorem]{Remark}
\newcommand{\Z}{\mathbb Z}
\newcommand{\C}{\mathbb C}
\newcommand{\R}{\mathbb R}
\newcommand{\Q}{\mathbb Q}
\newcommand{\Ann}{\operatorname{Ann}}
\newcommand{\supp}{\operatorname{supp}}
\newcommand{\Newt}{\operatorname{Newt}}
\newcommand{\Pat}{\operatorname{Pat}}
\newcommand{\conv}{\operatorname{conv}}
\newcommand{\rec}{\operatorname{rec}}
\newcommand{\Span}{\operatorname{span}}
\newcommand{\Laur}{\mathscr L}
\newcommand{\Del}[1]{\Delta_{#1}}
\newcommand{\Rect}{R_{k,n}}
\newcommand{\ZZ}{\mathcal Z}
\newcommand{\RR}{\mathcal R}
\DeclareMathOperator{\lcm}{lcm}
\title{A proof of Nivat's conjecture}
\author{Pranava Kumar}
\date{14 September 2026}
\begin{document}
\maketitle
\begin{abstract}
We prove that a finite-alphabet configuration on $\Z^2$ with at most $kn$ patterns on a $k\times n$
rectangle has a nonzero period. The main ingredient is a sharp lower bound: a nonperiodic
finite-valued configuration that has a product of directional binomial annihilators and agrees with a
fully periodic configuration on a rational wedge has at least $kn+1$ patterns on every $k\times n$
rectangle. We obtain this bound by supplementing
the affine coordinate functions with local quadratic functions. Their independence follows from a
support argument: directional differences produce nonzero pointwise products of half-plane-supported defects,
confined to pointed cones on which no Laurent recurrence can hold. A selection argument based on
Colle's construction supplies the wedge. Here we retain a finite boundary disagreement, so that the
orbit-closure point to which the bound is applied remains nonperiodic.
\end{abstract}

\section{Introduction}
Let $c:\Z^2\to\mathcal A$ be a configuration over a finite alphabet. For a finite set $D\subset\Z^2$, write
\[
 \Pat_c(D)=\{(c(g+u))_{u\in D}:g\in\Z^2\},\qquad P_c(D)=|\Pat_c(D)|.
\]
Thus $P_c(D)$ counts the patterns seen by translating the observation window $D$. We use the rectangles
\[
 \Rect=\{0,\ldots,k-1\}\times\{0,\ldots,n-1\},\qquad
 P_c(k,n)=P_c(\Rect).
\]
A \emph{period} of $c$ is a nonzero vector $h\in\Z^2$ such that $c(g+h)=c(g)$ for all $g$. We call
$c$ \emph{fully periodic} if it has two linearly independent periods, and \emph{nonperiodic} if it
has no period. The following is the classical Nivat statement, a two-dimensional analogue of the Morse--Hedlund theorem \cite{MH}.

\begin{theorem}\label{thm:nivat}
Let $\mathcal A$ be a finite alphabet and $c:\Z^2\to\mathcal A$. If
\[
 P_c(k,n)\le kn
\]
for some positive integers $k,n$, then $c$ has a period.
\end{theorem}

The threshold is sharp. A configuration equal to $1$ at the origin and $0$ elsewhere is nonperiodic
and has exactly $kn+1$ patterns on $\Rect$. The conclusion cannot be strengthened to full
periodicity: if $s$ is a finite-valued aperiodic sequence, then $c(x,y)=s(x)$ has a vertical period
but no transverse period, while $P_c(1,n)$ is at most the alphabet size and hence at most $n$ for
sufficiently large $n$.

Cyr and Kra proved periodicity under the bound $P_c(k,n)\le kn/2$
\cite[Theorem~1.2]{CK}. Their generating-set method was refined by Colle and Garibaldi to retain
information from the alphabet size \cite[Theorem~2.2]{CG}. Short rectangles admit the full
threshold: Cyr and Kra treat short side at most three \cite[Theorem~1.1]{CKshort}, and Colle and
Garibaldi treat short side at most four \cite[Theorem~1.3]{CG26}.

A different approach, due to Kari and Szabados, associates a Laurent annihilator to a low-complexity
configuration and obtains a product of directional differences \cite[Lemma~5 and Corollary~13]{KS}.
Szabados proves the conjecture when the resulting periodic decomposition has at most two components
\cite{Sz}. Kari and Moutot show that the orbit closure contains a periodic point, and
obtain the conjecture for uniformly recurrent configurations \cite[Theorem~5 and Corollary~8]{KM}.
These conclusions do not remove the need to control the particular point to which a complexity
estimate is applied. Our selection argument follows the periodic-region constructions of
Colle \cite[Sections~3--4]{Colle}; their relation to the work of Colle and Garibaldi is described in Appendix~\ref{app:sources}.

Even the single-defect example above has the zero configuration in its orbit closure: one can simply
translate the defect to infinity. Thus a periodic limit, or agreement with a periodic background on
a wedge, is not itself a contradiction. We need a point that retains nonperiodicity. On this point,
the periodic region will be used to produce additional patterns, rather than to extend a period
through the plane.

We use the shifts and differences
\[
 T^u c(g)=c(g+u),\qquad \Del u=T^u-1,
\]
and write $X_c=\overline{\{T^u c:u\in\Z^2\}}$ for the orbit closure in the product topology, with
the finite alphabet discrete. A complex configuration is \emph{finitary} if it has finite image. A
\emph{rational wedge} is a translate of a closed rational cone with nonempty interior and aperture
less than $\pi$.

\begin{theorem}[The periodic-wedge bound]\label{thm:wedge}
Let $c:\Z^2\to\C$ be finitary and nonperiodic. Suppose that
\[
 Q_0c=0,\qquad Q_0=\prod_{\ell=1}^q\Del{h_\ell},
\]
where the nonzero integer vectors $h_\ell$ have pairwise distinct unoriented directions. If $c$
agrees with a fully periodic configuration on a rational wedge, then
\[
 P_c(k,n)\ge kn+1\qquad(k,n\ge1).
\]
\end{theorem}

If the constant function and the $kn$ coordinate functions were independent, they would give the desired bound. Annihilators can impose relations among them, so we look for nonlinear functions of the
same pattern. Let $\phi$ be the common directional divisor of the annihilators. Applying all but one
of its directional factors gives a singly periodic configuration whose agreement with the background
spreads from the wedge to a half-plane. A transverse difference then vanishes on that half-plane,
after a fixed inset.

With two distinct directions, we take the \emph{pointwise} product of the resulting defects:
for arrays $x,y:\Z^2\to\C$, $(xy)(g)=x(g)y(g)$. Its support is the intersection of
the two supports, hence is contained in a translated pointed cone. This rules out a nonzero Laurent
annihilator provided the product itself is nonzero.
There is a placement issue here: shifting arbitrary nonzero sites together could require shifts far
larger than the rectangle. Short recurrences along the factor directions instead allow us to work
inside the coordinate boxes of the factors of $\phi$. The resulting product is the double difference
of a quadratic local observable, whose fitting translates are independent modulo the affine
coordinate functions.

The same factor boxes also control how many affine relations there can be. Write $r,s$ for the
coordinate widths of $\phi$ and $D_R$ for the dimension of the affine relations on $R=\Rect$. With
$t_+=\max\{t,0\}$, we have
\[
 D_R\le L_R:=(k-r)_+(n-s)_+.
\]
There are exactly $L_R$ translates of the quadratic observable that fit in $R$. Adding them to the
$kn+1-D_R$ dimensions supplied by the affine coordinate functions gives
\[
 P_c(k,n)\ge (kn+1-D_R)+L_R\ge kn+1.
\]
The details occupy Sections~\ref{sec:algebra}--\ref{sec:count}.

\begin{proposition}[Selection of a nonperiodic point]\label{prop:selection}
Suppose that $c:\Z^2\to\Z$ is finitary and nonperiodic and $P_c(k,n)\le kn$. There exist a
nonperiodic $\theta$ in the orbit closure $X_c$, a fully periodic configuration $B$, and a rational
wedge $K$ such that $\theta=B$ on $K\cap\Z^2$.
\end{proposition}

We prove this in Sections~\ref{sec:selection-inputs}--\ref{sec:selection-proof}, adapting the two
cases in the proof of Colle's Lemma~4.6 \cite{Colle}. The main difficulty in the compactness case is to retain nonperiodicity. We anchor
maximal agreement polygons at a boundary point and keep a disagreement on a fixed finite set. A
hypothetical period of the limit would then extend agreement to a reference half-plane, or give that
half-plane a second period; either conclusion is ruled out by the construction. The agreement region
spreads to that half-plane under the reference period. No recurrence hypothesis is needed.

\begin{proof}[Deduction of Theorem~\ref{thm:nivat}]
Relabel the alphabet injectively by integers, preserving patterns and periods, and suppose for a
contradiction that $c$ is nonperiodic. If $n=1$, every horizontal row has at most $k$ words of
length $k$. The Morse--Hedlund theorem \cite[Theorem~2.2(iii)]{CK} gives each row a period at most $k$, so $\lcm(1,\ldots,k)$ is a common horizontal period. The case $k=1$ follows in the same way from the columns.

We may therefore assume $k,n\ge2$. Lemma~\ref{lem:affine-ann} supplies a nonzero integral annihilator, and Theorem~\ref{thm:KS} supplies the product $Q_0$. This finite local identity holds on every point of $X_c$. Every finite pattern of such a point also occurs in $c$, since its cylinder neighborhood meets the orbit. Apply Proposition~\ref{prop:selection} and then Theorem~\ref{thm:wedge} to the selected nonperiodic point $\theta$. We obtain
\[
 kn+1\le P_\theta(k,n)\le P_c(k,n)\le kn,
\]
a contradiction.
\end{proof}

\section{Annihilators and directional filters}\label{sec:algebra}
We regard configurations as functions on $\Z^2$, rather than multiply them as formal Laurent
series. Their values may be added and multiplied pointwise. In particular,
\[
 (xy)(g)=x(g)y(g),\qquad
 \supp(x):=\{g\in\Z^2:x(g)\ne0\}.
\]
No summation is involved, even when both supports are infinite. Over $\C$,
\begin{equation}\label{eq:pointwise-support}
 \supp(xy)=\supp(x)\cap\supp(y),\qquad
 \supp(T^u x)=\supp(x)-u.
\end{equation}
The second sign follows from our convention $T^u x(g)=x(g+u)$.

Laurent polynomials play a different role. Put $\Laur=\C[X^{\pm1},Y^{\pm1}]$ and abbreviate
$X^{h_1}Y^{h_2}$ by $X^h$. For $f=\sum_{u\in\supp(f)}f_uX^u\in\Laur$, define
\begin{equation}\label{eq:finite-action}
 f(T)x(g)=\sum_{u\in\supp(f)} f_u x(g+u),\qquad
 \Ann(x)=\{f\in\Laur:f(T)x=0\}.
\end{equation}
This sum is finite. Polynomial multiplication corresponds to composition,
$(fg)(T)x=f(T)(g(T)x)$, so $\Ann(x)$ is an ideal. This is the annihilator formalism of
Kari--Szabados \cite[Section~2]{KS}, with the shift convention reversed. In formal-series notation,
our product of two configurations would be a coefficientwise (Hadamard) product, not a Cauchy
product. No convolution of two infinitely supported arrays is used.

A shift preserves pointwise products: $T^u(xy)=(T^u x)(T^u y)$. A general operator $f(T)$ does
not. We will expand the two differences of a product explicitly when this distinction matters.
Here, \emph{finitary} means finite image, not finite support. Finite sums and pointwise products
of translates of a finitary configuration remain finitary.

For a polynomial $f$, the set $\supp(f)$ is the input stencil of $f(T)$; it need not contain the
support of the output array $f(T)x$. Its Newton polygon is
$\Newt(f)=\conv(\supp(f))$ for $f\ne0$. The coordinate minima and maxima of nonzero Laurent
polynomials add under multiplication: viewing a
polynomial in $X$ with coefficients in the integral domain $\C[Y^{\pm1}]$ proves this for the two
$X$-extrema, and interchanging the coordinates proves it for the other two. The \emph{coordinate
widths} are the differences between the corresponding maxima and minima. We call a nonconstant
Laurent polynomial a \emph{line polynomial} when its Newton polygon is a nondegenerate segment.

The initial linear-algebra argument is the one in Kari--Szabados \cite[Lemma~5]{KS}. We include it
to keep the equality case visible.

\begin{lemma}\label{lem:affine-ann}
If $c$ is finitary and integer-valued and $P_c(D)\le |D|$ for a nonempty finite set $D$, then $c$
has a nonzero integral Laurent annihilator.
\end{lemma}
\begin{proof}
Fix a $D$-pattern $p_0$. The at most $|D|-1$ nonzero vectors $p-p_0$ have a nonzero orthogonal
vector in $\Q^D$. Its coefficients define a nonzero polynomial $f$ supported in $D$ for which
$f(T)c$ is constant. For any $w\in\Z^2\setminus\{0\}$, the nonzero polynomial $(X^w-1)f$ annihilates
$c$. Clearing denominators gives an integral annihilator.
\end{proof}

\begin{theorem}[Kari--Szabados, Corollary 13 of \cite{KS}]\label{thm:KS}
A finitary integer configuration with a nonzero annihilator has an annihilator
$\prod_{\ell=1}^q(X^{h_\ell}-1)$ with nonzero integer vectors in pairwise distinct directions.
\end{theorem}
The source uses the opposite shift convention. Replacing every exponent by its negative gives the
statement above, since reversing a period binomial changes it only by a Laurent unit. We use the
product identity, not a bounded decomposition into periodic components.

The next proof uses the elimination argument of Kari--Moutot \cite[Lemma~8]{KMalg}. We state the
finite-image hypothesis explicitly: it is what turns the resulting recurrences into periods.

\Needspace{7\baselineskip}
\begin{lemma}[Coprime annihilators]\label{lem:coprime}
A finitary complex configuration with two coprime nonzero Laurent annihilators is fully periodic.
\end{lemma}
\begin{proof}
Let $f,g$ be the two annihilators. They are coprime in $\C(X)[Y^{\pm1}]$: a common nonunit factor
there would give one in $\Laur$ after clearing denominators and taking primitive parts. A B\'ezout
identity in this one-variable ring, with denominators cleared, gives a nonzero polynomial in $X$ in
$(f,g)$. Interchanging the variables gives a nonzero polynomial in $Y$.

We remove the monomial factor from each eliminant. A nonconstant eliminant now gives a reversible
recurrence of some order $d$ along every parallel lattice line, since its constant and leading
coefficients are nonzero. The length-$d$ states occurring anywhere in the configuration form one
finite set. Each state determines both its successor and its predecessor, so the successor map is a
permutation, whose order is a period on all the parallel lines simultaneously. Applying this to the
two eliminants gives independent coordinate periods. If either eliminant is a nonzero constant,
the configuration is zero.
\end{proof}

For the remainder of Sections~\ref{sec:algebra}--\ref{sec:count}, assume the hypotheses of
Theorem~\ref{thm:wedge}. An individual annihilator may have factors which another annihilator lacks.
For the local count we need their common factors. We do not assume that their gcd annihilates $c$; the
following lemma shows that applying it leaves only a fully periodic remainder
(cf.~\cite[Corollaries~28--29]{KS}).

\begin{lemma}[The common directional divisor]\label{lem:gcd}
The nonzero elements of $\Ann(c)$ have a gcd $\phi$, defined up to a Laurent unit. It can be written
\[
 \phi=\phi_1\cdots\phi_m,\qquad m\ge2,
\]
where the $\phi_i$ are nonconstant line polynomials in distinct directions, each dividing the
corresponding binomial in $Q_0$. Choose their monomial representatives so that both coordinate
minima of every $\phi_i$ are zero, and use their product as the representative of $\phi$. Then
$\rho=\phi(T)c$ is fully periodic. In addition,
\[
 f(T)c\text{ constant}\quad\Longrightarrow\quad \phi\mid f.
\]
\end{lemma}
\begin{proof}
Write $h=ap$ with $a\ge1$ and $p$ primitive. Then the factors of $X^h-1$ over $\C$ are $X^p-\zeta$,
with $\zeta^a=1$, up to a unit. They are distinct. Binomials in different directions have no common
nonunit factor. Thus $Q_0$ is squarefree.

Because $Q_0\in\Ann(c)$, we may define $\phi$ as the product of the irreducible factors of $Q_0$
which divide every annihilator. This is a gcd even though the ideal contains infinitely many
polynomials. For each factor omitted from $\phi$, choose an annihilator not divisible by that
factor. We can choose a finite linear combination $g$ of these annihilators which avoids every
omitted factor: each divisibility condition defines a proper linear subspace of the coefficient
space, and a finite union of such subspaces does not cover a complex vector space. Thus
$\gcd(Q_0,g)=\phi$. If no factor is omitted, take $g=Q_0$.

The coprime quotients $Q_0/\phi$ and $g/\phi$ annihilate $\rho=\phi(T)c$. This configuration is
finitary, so Lemma~\ref{lem:coprime} makes it fully periodic. Grouping the irreducible factors of
$\phi$ by direction, we obtain the polynomials $\phi_i$. Multiplying each by the prescribed monomial unit
fixes its representative without changing the divisibility statements or periodicity.

There cannot be zero or one surviving direction. The first case makes $c$ fully periodic. In the
second, choose a nonzero multiple $v$ of the vector in $Q_0$ which is also a period of $\rho$, and
write $\Del v=\phi a$. The difference $\Del v c=a(T)\rho$ is $v$-periodic, so its increments along
each $v$-line are constant: for $N\ge0$,
\[
 c(g+Nv)=c(g)+N(a(T)\rho)(g).
\]
Boundedness forces $a(T)\rho=0$, making $v$ a period of $c$. Both cases contradict nonperiodicity, so $m\ge2$.

If $f(T)c$ is constant, choose $w\in\Z^2\setminus\{0\}$ parallel to none of the directions of $\phi$. Then $\Del w
f\in\Ann(c)$ and $\gcd(\Del w,\phi)=1$. Every irreducible factor of the squarefree polynomial $\phi$
therefore divides $f$, as required.
\end{proof}

Put
\[
 f_i=\phi/\phi_i,\qquad A_i=f_i(T)c.
\]
The filter $f_i$ applies every directional factor except $\phi_i$. The phrase ``direction'' refers
to a line, not a chosen period: a primitive vector spans all lattice points on that line, but need
not itself be a period. Let $B$ be the fully periodic background on the wedge. Choose a nonzero multiple $v_i$ of the corresponding vector in $Q_0$ which
is a period of both $\rho$ and $B$. This is possible because both period lattices have finite index
in $\Z^2$.

\begin{lemma}[Directional filters]\label{lem:filters}
Each $A_i$ has period $v_i$, is not fully periodic, and agrees with a fully periodic configuration
on a rational lattice half-plane whose boundary is parallel to $v_i$.
\end{lemma}
\begin{proof}
Write $\Del{v_i}=\phi_i a_i$. Then
\[
 \Del{v_i}A_i=a_i(T)\rho
\]
is $v_i$-periodic. Applying the boundedness argument of Lemma~\ref{lem:gcd} to $A_i$ gives
$\Del{v_i}A_i=0$. A transverse period $w$ would put $\Del w f_i$ in $\Ann(c)$, although it is
coprime to $\phi_i$. Thus $A_i$ is not fully periodic.

A finite stencil only costs a translation of the wedge into its interior. After making this
translation, $A_i=B_i$ on a smaller wedge $K'$, where $B_i=f_i(T)B$. Their common period $v_i$
spreads agreement to $(K'\cap\Z^2)+\Z v_i$. For a nonprimitive period this set can miss lattice
points near its boundary; what we need, and will prove, is that it contains an inset lattice
half-plane.

Write $K'=\{g:\lambda_1(g)\ge a_1,\ \lambda_2(g)\ge a_2\}$ with independent rational linear forms. A
lattice point $g$ belongs to this saturation whenever some integer $t$ satisfies
\[
 \lambda_1(g)+t\lambda_1(v_i)\ge a_1,\qquad
 \lambda_2(g)+t\lambda_2(v_i)\ge a_2.
\]
If $\lambda_1(v_i)$ and $\lambda_2(v_i)$ have the same strict sign, an integer $t$ exists for every
$g$. This would give $A_i=B_i$ globally, contrary to non-full-periodicity. If one coefficient
vanishes, the other is nonzero, and the inequalities can be solved throughout the half-plane imposed
by the inequality with zero coefficient of $t$.

In the remaining case, after reversing $v_i$ if necessary, write $a=\lambda_1(v_i)>0$ and
$b=-\lambda_2(v_i)>0$. The admissible interval is
\[
 \frac{a_1-\lambda_1(g)}{a}\le t\le
 \frac{\lambda_2(g)-a_2}{b}.
\]
Put $\nu=\lambda_1/a+\lambda_2/b$. This is a nonzero rational form with $\nu(v_i)=0$. On the half-plane
\[
 \nu(g)\ge \frac{a_1}{a}+\frac{a_2}{b}+1,
\]
the interval has length at least $1$ and contains an integer. This proves the assertion for
nonprimitive periods as well; the inset avoids any missing lattice points near the boundary of the
real saturation.
\end{proof}

\section{The pointed-cone obstruction and a local observable}\label{sec:observable}
An array supported in a half-plane may still be periodic along its boundary. A pointwise product
of two such arrays has a different support restriction: by \eqref{eq:pointwise-support}, it vanishes
outside the intersection of their half-planes. When the boundaries are nonparallel this intersection
is a translated pointed cone. The product need not be nonzero; Lemma~\ref{lem:placement} will deal
with that issue. We first record what the support restriction gives once nonvanishing is known.

\begin{lemma}[The pointed-cone obstruction]\label{lem:cone}
A nonzero complex array supported in a translate of a closed pointed cone of aperture less than
$\pi$ has no nonzero Laurent annihilator.
\end{lemma}
\begin{proof}
Choose a linear functional $\lambda$ in the interior of the dual cone whose kernel contains no
nonzero lattice vector; this requires avoiding only countably many lines. Its sublevel sets in the
translated cone are bounded and contain finitely many lattice points. Hence the nonempty support of
the array has a unique $\lambda$-minimum $g_*$. For a nonzero operator $f=\sum_u a_uT^u$, choose the
unique $\lambda$-maximum $u_*$ of its exponent support. In
\[
 f(T)e(g_*-u_*)=\sum_u a_u e(g_*-u_*+u),
\]
all terms with $u\ne u_*$ sample points below the minimum $g_*$ and hence vanish. The remaining
term is $a_{u_*}e(g_*)\ne0$. We have used a maximum for the operator and a minimum for the array
because $T^u$ samples at $g+u$.
\end{proof}

Let $r_i,s_i$ be the coordinate widths of the normalized $\phi_i$, and put $r=\sum_i r_i$, $s=\sum_i
s_i$. Let $[E_i^-,E_i^+]$ be its Newton segment, of primitive lattice length $t_i\ge1$. The support
of $f_i$ lies in
\[
 [0,r-r_i]\times[0,s-s_i].
\]
Fix distinct indices $i,j$. Increase $v_i,v_j$ by common multiples so that they fix $\rho$ and all
the fully periodic $B_i$ in Lemma~\ref{lem:filters}. Define the transverse defects
\[
 b_i=\Del{v_j}A_i,\qquad b_j=\Del{v_i}A_j.
\]
Neither is zero, since otherwise the corresponding filtered configuration would be fully periodic.
To locate the support, write an agreement half-plane for $A_i,B_i$ as $\nu_i(g)\ge\beta_i$, with
$\nu_i(v_i)=0$. Since $v_j$ is a period of $B_i$,
\[
 b_i(g)=0\qquad\text{if}\qquad
 \nu_i(g)\ge\beta_i+\max\{0,-\nu_i(v_j)\}.
\]
Thus $b_i$ is supported in the opposite closed half-plane. The same argument applies to $b_j$. Their
support boundaries are parallel to $v_i,v_j$, respectively. Also
\begin{equation}\label{eq:small-recurrence}
 \phi_i(T)b_i=\Del{v_j}\rho=0,\qquad
 \phi_j(T)b_j=\Del{v_i}\rho=0.
\end{equation}
The useful length scale is that of these recurrences, not the possibly large periods $v_i,v_j$.
On a line where a defect is nonzero, a short recurrence prevents a long run of zeros. We can
therefore move nonzero values together using only the factor boxes.

\begin{lemma}[Placement in the factor boxes]\label{lem:placement}
There are integer vectors $u_i\in[0,r_i]\times[0,s_i]$ and $u_j\in[0,r_j]\times[0,s_j]$ for which
\[
 e=(T^{u_i}b_i)(T^{u_j}b_j)
\]
is nonzero.
\end{lemma}
\begin{proof}
Choose a real affine line $L_i$ in the direction of $\phi_i$ containing a lattice point where $b_i$
is nonzero, and choose $L_j$ in the same way. The translated lines intersect at
\[
 z=(L_i-E_i^-)\cap(L_j-E_j^-).
\]
Although $z$ need not be integral, the segment $z+[E_i^-,E_i^+]$ lies on $L_i$ and contains at least
$t_i$ consecutive points of $L_i\cap\Z^2$: an interval of length $t_i$ in primitive lattice units
contains at least $t_i$ lattice points, regardless of its endpoints.

Translate the operator by $-E_i^-$. Equation~\eqref{eq:small-recurrence} then gives a recurrence of
order $t_i$ on $L_i\cap\Z^2$ with nonzero first and last coefficients. A run of $t_i$ zeros would
propagate in both directions and make the entire line zero. We may therefore choose $p_i$ on the
segment with $b_i(p_i)\ne0$, and similarly choose $p_j$ for $b_j$.

Set $g=\lfloor z\rfloor$ coordinatewise and $u_i=p_i-g$, $u_j=p_j-g$. Since $p_i-z$ belongs to the
integral coordinate box $[0,r_i]\times[0,s_i]$,
\[
 u_i=p_i-\lfloor z\rfloor=\lceil p_i-z\rceil
\]
also belongs to that box. The same holds for $j$, and $e(g)=b_i(p_i)b_j(p_j)\ne0$.
\end{proof}

If $H_i,H_j$ are the half-planes containing the supports of $b_i,b_j$, then
\[
 \supp(e)=(\supp(b_i)-u_i)\cap(\supp(b_j)-u_j)
       \subset (H_i-u_i)\cap(H_j-u_j).
\]
The containing half-planes, not the supports themselves, intersect in a translated pointed cone.
Since $e\ne0$, Lemma~\ref{lem:cone} applies. We do not use $e$ directly as the local observable,
since its definition involves the large shifts $v_i,v_j$. Instead, form the pointwise product before
taking those differences:
\[
 d=(T^{u_i}A_i)(T^{u_j}A_j).
\]
Its input sites, relative to the evaluation site, belong to
\begin{equation}\label{eq:observable-stencil}
 (u_i+\supp(f_i))\cup(u_j+\supp(f_j))
       \subset W:=[0,r]\times[0,s]\cap\Z^2.
\end{equation}
Thus $d(g)$ is a quadratic function of the $c$-values on $g+W$. This is a statement about its
input stencil, not a claim that the output array $d$ has finite support.

To recover $e$, write $U(g)=A_i(g+u_i)$ and $V(g)=A_j(g+u_j)$, so that $d(g)=U(g)V(g)$.
Since $T^{v_i}U=U$ and $T^{v_j}V=V$,
\begin{align*}
 \Del{v_i}\Del{v_j}(UV)
 &= (T^{v_j}U)(T^{v_i}V)-U(T^{v_i}V)-(T^{v_j}U)V+UV\\
 &= (\Del{v_j}U)(\Del{v_i}V)=e.
\end{align*}
Set $Q=\Del{v_i}\Del{v_j}Q_0$. This operator annihilates both $c$ and constants, while
\begin{equation}\label{eq:Qd}
 Qd=Q_0e
\end{equation}
has no nonzero annihilator. Indeed, an annihilator $H$ of $Q_0e$ would make the nonzero polynomial
$HQ_0$ annihilate $e$.

\section{The exact rectangle count}\label{sec:count}
Kari--Szabados relate linear independence of patterns to the absence of annihilators fitting in a
window \cite[Lemma~36]{KS}. We keep the constant coordinate as well, rather than first normalizing
the configuration as in \cite[Lemma~33]{KS}. The extra step is to supplement this affine space
with translates of the quadratic function just constructed.

\begin{proof}[Proof of Theorem~\ref{thm:wedge}]
Fix $R=\Rect$ and let
\[
 F_R=\{f\in\Laur:\supp(f)\subset R,\ f(T)c\text{ is constant}\},\qquad D_R=\dim F_R.
\]
Lemma~\ref{lem:gcd} makes every member of $F_R$ divisible by $\phi$. Since the latter has coordinate
minima zero and widths $r,s$, additivity of coordinate extrema confines every nonzero quotient
$f/\phi$ to
\[
 [0,k-r-1]\times[0,n-s-1].
\]
Multiplication by $\phi$ is injective, so
\begin{equation}\label{eq:relation-count}
 D_R\le L_R:=(k-r)_+(n-s)_+.
\end{equation}
If $k\le r$ or $n\le s$, the quotient box has no lattice sites and $D_R=0$.

Let $V_R$ be the subspace of all complex functions on $\Pat_c(R)$ spanned by $1$ and the coordinate
evaluations $p\mapsto p(u)$, $u\in R$. The map from their $kn+1$ coefficients to this function space
has kernel
\[
 \{(-\kappa(f),f):f\in F_R\},\qquad f(T)c=\kappa(f).
\]
Since $\kappa(f)$ is uniquely determined by $f$, including for $f=0$, we obtain
\begin{equation}\label{eq:affine-count}
 \dim V_R=kn+1-D_R.
\end{equation}
If $L_R=0$, this already proves the theorem. For integral configurations, this thin-rectangle
conclusion also appears in Kari--Szabados \cite[Lemma~39]{KS}.

Suppose $L_R>0$. Write $f_i=\sum_a\alpha_{i,a}X^a$ and
$f_j=\sum_b\alpha_{j,b}X^b$, with both sums restricted to the finite polynomial supports. The
function on $W$-patterns which gives $d$ is
\[
 \mathcal D(p)=
 \left(\sum_a\alpha_{i,a}p(u_i+a)\right)
 \left(\sum_b\alpha_{j,b}p(u_j+b)\right).
\]
In other words, $d(g)=\mathcal D((c(g+z))_{z\in W})$. This also makes the word ``local'' precise.
For each $t$ with $t+W\subset R$, the function
$p\mapsto\mathcal D((p(t+z))_{z\in W})$ is defined on $\Pat_c(R)$; evaluating it at the pattern
based at $g$ gives $T^td(g)$. There are exactly $L_R$ such translations.

We claim that these functions are independent modulo $V_R$. A relation on the finite set
$\Pat_c(R)$, evaluated at every translate of $c$, gives the global array identity
\[
 \sum_t\alpha_t T^td=\alpha+f(T)c,\qquad \supp(f)\subset R.
\]
The identity holds at every $g\in\Z^2$, so it remains valid at all sites sampled by $Q$.
Applying $Q$ removes the affine terms and gives
\[
 \Big(\sum_t\alpha_tT^t\Big)(Qd)=0.
\]
By \eqref{eq:Qd}, all $\alpha_t$ vanish. These $L_R$ functions are thus independent modulo $V_R$.
The full function space on the finite set $\Pat_c(R)$ has dimension $P_c(R)$, and consequently
\[
 P_c(k,n)\ge\dim V_R+L_R
 =kn+1-D_R+L_R\ge kn+1.
\]
The distinction between defining the functions and testing a relation among them is worth keeping
in view. Each function counted above uses only the $R$-pattern; $Q$ acts afterwards on the global
identity, so its large shifts do not enlarge the observation window.
\end{proof}

\section{Generating sets and a periodic reference}\label{sec:selection-inputs}
For Sections~\ref{sec:selection-inputs}--\ref{sec:selection-proof}, fix a nonperiodic finitary
integer configuration $c$ with $P_c(k,n)\le kn$ and a product annihilator $Q_0$ as in
Theorem~\ref{thm:KS}. Since $c$ is nonperiodic, $Q_0$ has at least two factor directions.
Every point of $X_c$ satisfies $Q_0z=0$ and uses only finite patterns occurring in $c$. The reference we construct below will have one period, but its restriction to a chosen
half-plane must not be fully periodic. This distinction will later exclude a transverse period of
the selected point.

A finite set is \emph{lattice-convex} if it consists of all lattice points of its convex hull. A
vertex $v$ of a lattice-convex set $S$ is \emph{generated} if any two occurring $S$-patterns
agreeing on $S\setminus\{v\}$ also agree at $v$. We call $S$ a \emph{generating set} when every
vertex is generated, and an array \emph{$S$-legal} when each of its translated $S$-patterns occurs
in $c$. An edge $e$ in a pattern count denotes its lattice sites; $\conv(e)$ denotes the
corresponding geometric segment. We use the generating-set construction of Cyr--Kra
\cite[Lemmas~2.5 and~2.8]{CK}, in the extremal form stated by Colle \cite[Section~2]{Colle}.
The short discrepancy proof below records the edge bound needed later.

\begin{lemma}[A generating polygon]\label{lem:generating}
There is a lattice-convex $S\subset\Rect$ of positive area such that every vertex is generated and,
for every edge $e$ of $S$,
\begin{equation}\label{eq:edge-budget}
 P_c(S)-P_c(S\setminus e)\le |e|-1.
\end{equation}
\end{lemma}
\begin{proof}
Minimize the discrepancy $P_c(S)-|S|$ over nonempty lattice-convex subsets of $\Rect$, and choose an
inclusion-minimal minimizer. Its discrepancy is nonpositive, since $\Rect$ is a competitor.

If $S$ were a segment of $t$ consecutive lattice sites, every parallel bi-infinite line would have
at most $t$ words of length $t$. Morse--Hedlund would give each line a period at most $t$, and their
common multiple would be a global period of $c$. This also excludes a singleton, so $S$ has positive
area.

Every proper nonempty lattice-convex subset of $S$ has larger discrepancy. Deleting a vertex $v$
therefore gives $P_c(S)-P_c(S\setminus\{v\})<1$. This difference is a nonnegative integer, since
restriction of patterns is surjective, so it is zero and $v$ is generated. Deleting an edge also
leaves a nonempty lattice-convex set, and the same comparison gives \eqref{eq:edge-budget}.
\end{proof}

For an edge $e$, an occurring $S\setminus e$-pattern is \emph{ambiguous} if it has at least two
occurring $S$-extensions. If $N(\gamma)$ is the number of such extensions, then
\begin{equation}\label{eq:ambiguity-budget}
 \sum_{\gamma}\bigl(N(\gamma)-1\bigr)
 =P_c(S)-P_c(S\setminus e)\le |e|-1.
\end{equation}
In particular, there are at most $|e|-1$ ambiguous base patterns.

Following Colle \cite[Section~1]{Colle}, an oriented line through the origin is
\emph{one-sided nonexpansive} for $X_c$ if two distinct points of $X_c$ agree on the closed lattice
half-plane on its left. The two orientations specify the two sides of the line. The next result is
the contrapositive of Colle's theorem.

\begin{theorem}[Colle, Theorem 1.9 of \cite{Colle}]\label{thm:paired}
If a finitary integer configuration with a nonzero annihilator is not fully periodic, then some line
has both orientations one-sided nonexpansive on its orbit closure.
\end{theorem}

Both supporting faces of the generating polygon in this direction must be edges. Indeed, suppose the
face were a single generated vertex $v$, and let $\lambda$ increase into the half-plane. The
positive minimum of $\lambda(u-v)$ over $u\in S\setminus\{v\}$ gives a slab thickness through which
agreement propagates. Iterating would determine the entire plane, contradicting one-sided
nonexpansiveness. In particular the line is rational. Applying the same argument to the nonzero
extreme coefficient of $Q_0$ shows that the line is parallel to a factor of $Q_0$: otherwise its
Newton polygon has a unique minimizing vertex.

A nonzero vector $h$ is a \emph{regional period} of $x$ on $U$ if $x(g+h)=x(g)$ whenever $g,g+h\in
U$. The restriction is \emph{fully periodic} if it has two linearly independent regional periods
(compare \cite[Definition~2.12]{Colle}). If $U+h\subset U$, a regional period $h$ is equivalently
expressed by $\Del h x=0$ on $U$. We use this equivalence only when the indicated translation
stays in the region.

For a tangent half-plane period, we can work directly with its difference. Taking a suitable
multiple of the period cancels any parallel factor of the annihilator; the remaining factors then
determine successive normal levels, extending the period to the plane
(cf.~\cite[Proposition~2.14]{Colle}).

\begin{lemma}[Extending a tangent period]\label{lem:globalize}
Suppose $Q_0z=0$ and $z$ has a period parallel to the boundary of a rational half-plane on that
half-plane. Then $z$ has a global period in the same direction.
\end{lemma}
\begin{proof}
If a factor $\Del{h_j}$ is parallel to the boundary, choose a nonzero common multiple $h$ of $h_j$
and the regional period, and write $\Del h=a\Del{h_j}$. Then
\[
 E=\Del h z,\qquad \Bigl(Q_0/\Del{h_j}\Bigr)E=aQ_0z=0.
\]
If there is no parallel factor, take the given period difference as $E$ and retain $Q_0$. In either
case $E$ vanishes on the half-plane and has a product annihilator whose directions are all
transverse to the boundary.

Choose a primitive integer normal $\nu$ pointing into the half-plane, and reverse the remaining
factor vectors so that their $\nu$-values are positive. This changes the product only by a Laurent
unit. Its constant coefficient is nonzero, while every other exponent has positive integer
$\nu$-value. The recurrence therefore propagates zero one outward lattice level at a time, giving
$E=0$ on the plane. If the remaining product is empty, the same conclusion is immediate.
\end{proof}

The following construction adapts the ambiguity argument of Cyr--Kra \cite[Lemma~2.24]{CK},
as developed in Colle \cite[Proposition~2.13]{Colle}.

\begin{lemma}[A periodic reference]\label{lem:reference}
There are a primitive integer vector $a$, an integer half-plane $H$ with boundary parallel to $a$,
and a globally $a$-direction-periodic $w\in X_c$ such that $w|_H$ is not fully periodic. The
generating polygon $S$ has two opposite edges parallel to $a$, and this direction occurs in $Q_0$.
\end{lemma}
\begin{proof}
Choose the line supplied by Theorem~\ref{thm:paired}, and orient its primitive vector $a$ so that
the facing edge $e$ of $S$ is no longer than the opposite edge. Set $r=|e|-1\ge1$. Subtracting $r$
primitive units from these two opposite edges leaves nonnegative cyclic edge vectors with sum zero.
Hence $[0,ra]$ is a Minkowski summand of $\conv S$, and every parallel lattice row through $S$
contains at least $r$ consecutive sites.

Choose $x,y\in X_c$ agreeing on the associated half-plane and differing on the first outer lattice
line where agreement fails; translate the pair so that this is the next line. The first such line
exists because the normal levels are integral. Place $e$ on it. Every $a$-translate of the common
$S\setminus e$-pattern is ambiguous: a unique extension at one position would give agreement on an
edge block, and generation of the two endpoints would extend this agreement along the entire line.

Let $d$ be the width of $S$ in primitive normal levels. On each of the $d$ rows of $(S\setminus
e)+\Z a$, choose the initial site of a run of $r$ sites in $S\setminus e$. Read these initial sites
as one tuple and translate the tuple along $a$. Every length-$r$ word of the resulting bi-infinite
sequence is a restriction of an ambiguous base pattern. There are at most $r$ such patterns by
\eqref{eq:ambiguity-budget}, so Morse--Hedlund gives the whole strip a common period $pa$, with
$p\le r$.

The period can change as we add an outer row. It is enough to keep its length at most $2r$ on each
strip of $d$ rows, since we can take a common multiple afterwards. A unique next-line extension
among $S$-legal configurations preserves the current period. Otherwise every translated base is
ambiguous by the endpoint argument, and the tuple argument gives a period $p_0a$ with $p_0\le r$.
The number of full $S$-patterns extending the observed bases is then at most
\[
 \sum_{\gamma\text{ observed}}N(\gamma)
 \le r+\#\{\gamma\text{ observed}\}\le2r.
\]
Among $2r+1$ successive full $S$-patterns of $x$, two therefore agree at a gap
$q\in\{1,\ldots,2r\}$. The agreement contains $r$ consecutive sites on every old row. Since these
rows are $p_0a$-periodic and $p_0\le r$, $qa$ is a period on the old strip. Generation of the
endpoints of $e$ propagates the equality in both directions along the new line, so the extended
strip has period $qa$.

Iterating gives each row in the outward half-plane a period of length at most $2r$. Thus
$\lcm(1,\ldots,2r)a$ is a common period there. The argument applies to both $x$ and $y$, and
Lemma~\ref{lem:globalize} makes each globally periodic in direction $a$.

Let $H$ be the original agreement half-plane enlarged to include the line where $x$ and $y$ differ.
At least one of $x|_H,y|_H$ is not fully periodic. To see this, a fully periodic restriction on a
half-plane extends to a fully periodic array: choose an inward regional period and transport each
point far enough into the half-plane. The values stabilize, and both regional periods become global.
If both restrictions had such extensions, the extensions would agree on the smaller half-plane. An
inward vector common to their two finite-index period lattices would then give agreement on the
plane, contradicting the differing boundary line. Choose $w$ to be a member whose restriction to $H$
is not fully periodic, and retain this choice throughout the selection argument.
\end{proof}

\section{Boundary continuation in polygons}\label{sec:geometry}
Let $\ZZ=\Newt(Q_0)$. Then $\ZZ$ is the Minkowski sum of the segments $[0,h_\ell]$, hence a lattice
zonotope. Interior subset sums may cancel in $Q_0$; the Newton polygon identity is unaffected.
Reversing a factor translates $\ZZ$.

A finite lattice polygon is \emph{$\ZZ$-enveloped} if it has the same oriented edge directions as
$\ZZ$, with each primitive edge length at least the corresponding length of $\ZZ$. Subtracting the
cyclic edge-vector lists writes it as $\ZZ+C$ for a compact convex polygon $C$, possibly a segment
or a point. A \emph{two-ray region} is a rational convex polyhedron with pointed two-dimensional
recession cone; its boundary consists of two rays joined by a finite chain of edges. In the
continuation arguments below, these regions have the form $\ZZ+C$ with $C$ a rational polyhedron.
Here $A+B=\{a+b:a\in A,\ b\in B\}$ denotes a Minkowski sum, and
$\rec(C)=\{v:C+tv\subset C\text{ for all }t\ge0\}$ is the recession cone. A facet is an exposed
one-dimensional face, which may be a finite edge or an infinite ray. Restrictions of arrays to real
regions always mean restrictions to their lattice points. The enveloped regions follow Colle's
geometric construction \cite[Definition~3.3]{Colle}; we retain the Minkowski summand because
it guarantees that the recurrence stencil fits at the boundary.

Once the interior values vanish, the recurrence on a new layer involves only an exposed face of
the Newton polygon. It becomes a one-dimensional recurrence along that layer. We will use this to control a growing boundary from a fixed finite seed.

\Needspace{16\baselineskip}
\begin{lemma}[Continuation across a facet]\label{lem:facet}
Let $f$ be a nonzero Laurent polynomial with Newton polygon $\ZZ$, and let $f(T)\delta=0$. Suppose
$P^+=\ZZ+C$ has a facet on $\nu=\beta$, with primitive integer inward normal $\nu$ and integer
$\beta$, and $\delta=0$ on
\[
 P=P^+\cap\{\nu\ge\beta+1\}.
\]
Orient the facet by a primitive integer vector $b$, pointing towards infinity when the facet is a
ray. Write the corresponding exposed face of $\ZZ$ as $[q_0,q_1]$, with $q_1-q_0=tb$ and $t\ge0$.
The case $t=0$ means that the exposed face is a single vertex.

If the facet of $P^+$ is finite and $\delta$ vanishes on its first $t$ consecutive lattice points in
direction $b$, it vanishes on the entire facet. If the facet is a ray in direction $b$, vanishing on
a tail of the ray implies vanishing on the entire ray. When $t=0$, no boundary seed is needed in
either assertion.
\end{lemma}
\begin{proof}
Let $F_C$ be the corresponding exposed face of $C$. If $t=0$, additivity of faces under Minkowski
summation gives $x-q_0\in F_C$ for every facet point $x$. In the recurrence at $x-q_0$, only the
nonzero extreme coefficient times $\delta(x)$ lies on the face; all other inputs are in $P$.
Thus $\delta(x)=0$ without a seed. Assume henceforth that $t\ge1$.

On a finite facet, a point $x$ at least $t$ primitive units from its initial endpoint satisfies
$x-q_1\in F_C$. Thus
\[
 x-q_1+\ZZ\subset P^+.
\]
In the recurrence at $x-q_1$, every nonface input lies in $P$, its normal level having increased by
a positive integer. The other face inputs precede $x$ along the facet, so the nonzero coefficient at
$q_1$ determines $\delta(x)$. Starting with the first $t$ zeros proves the assertion. This also
covers rational endpoints: each subsequent lattice point is at least $t$ primitive units from the
initial endpoint.

For a ray, the exposed face of $P^+$ is $q_0+F_C$. The recurrence at $x-q_0$ is therefore valid for
every lattice point $x$ on the ray, with all nonface inputs in $P$. Its first and last face
coefficients are nonzero. An all-zero tail extends backwards along the ray by solving for the term
at $q_0$.
\end{proof}

To expand a region outwards, we relax a signed supporting inequality. For an infinite facet on the
integer level $\nu=\beta$ of a primitive inward normal, let $\RR^{(j)}$ be obtained by replacing
that inequality with $\nu\ge\beta-j$ and keeping all others. Thus
$\RR^{(0)}=\RR\subset\RR^{(1)}\subset\cdots$. Moving the infinite facet lengthens the neighboring
finite edge and retains the $\ZZ$ summand. For a finite polygon, any fixed number of such layers
also retains the summand once the moving edge is sufficiently long: its neighbors lengthen, while
its own length changes by a bounded amount.

\begin{lemma}[A failed extension is ambiguous]\label{lem:boundary-ambiguity}
Let $x,y\in X_c$ agree on a two-ray region $\RR$ with Minkowski summand $\ZZ$, but not on its next
outward lattice layer $\RR^{(1)}$ along the infinite edge in primitive direction $b$. Assume the
moving facet is on an integer level of its primitive inward normal. Then $S$ has a supporting edge
$e_b$ in that direction. After translating $S$ so that $e_b$ lies on the new layer, every
sufficiently far positive $b$-translate of the common $S\setminus e_b$-pattern is ambiguous.
\end{lemma}
\begin{proof}
Place the minimizing face of $S$ on the new lattice line. Every other irredundant inward normal of
$\RR$ is strictly positive on $b$, so after a sufficiently large positive $b$-translation all
remaining sites of $S$ lie in $\RR$. They remain there under further positive translations.

A singleton supporting face would give agreement on a tail of the new line by generation of that
vertex. For an edge, a unique extension of one translated base pattern would give agreement on an
edge block; generation of its final endpoint would then propagate agreement along the remaining
tail. In either case, the ray part of Lemma~\ref{lem:facet}, applied to $x-y$, would extend
agreement to the entire new layer. This contradiction proves that the face is an edge and every
sufficiently far base pattern is ambiguous.
\end{proof}

Two periods pointing along the recession rays of a two-ray region $\RR$ determine one fully periodic
background. If these periods are $h_1,h_2$, then $t=h_1+h_2$ lies in the interior of the recession
cone, and we may define
\begin{equation}\label{eq:background}
 B(g)=x(g+Nt)\quad\text{for all sufficiently large }N.
\end{equation}
The value stabilizes: an $h_1$ step followed by an $h_2$ step stays in the region and preserves the
value. For either $i$, move both $g$ and $g+h_i$ into $\RR$ by the same sufficiently large multiple
of $t$. The regional period gives $B(g+h_i)=B(g)$, so $B$ is fully periodic. For $g\in\RR$, the same
steps show $B(g)=x(g)$. Since $\RR$ contains a translate of its recession cone, agreement holds on a
rational wedge.

\section{From one regional period to two}\label{sec:two-periods}
Following Colle \cite[Lemma~4.2]{Colle}, we use the existing period and generated vertices to
extend equalities from an eventually periodic transverse half-strip into the widening region.
A downward sweep requires the same transverse period on arbitrarily high starting bands;
repetitions of one finite pattern provide this choice.

\begin{lemma}[A second regional period]\label{lem:regional}
Let $a,b$ be primitive integer vectors spanning a pointed cone of positive aperture, and let $\RR$
be a rational two-ray polyhedron with recession cone $\R_{\ge0}a+\R_{\ge0}b$. Suppose $x\in X_c$ has
a period $pa$, $p\ge1$, on $\RR$. Let $S$ be a lattice-convex generating set for the language of
$c$, with opposite edges parallel to $a$ and a supporting edge $e_b$ parallel to $b$. Assume
\[
 P_c(S)-P_c(S\setminus e_b)\le |e_b|-1,
\]
and that the inward normal of $e_b$ is strictly positive on $a$. Suppose $S$ can be placed so that
every sufficiently far $b$-translate of $S\setminus e_b$ lies in $\RR$ and is ambiguous. Then $x$
has two independent periods, parallel to $a,b$, on a smaller two-ray region with the same recession
cone.
\end{lemma}
\begin{proof}
Let $r\ge1$ be the smaller primitive length of the two $a$-parallel edges. As in
Lemma~\ref{lem:reference}, the segment $[0,ra]$ is a Minkowski summand of $\conv S$. Since the
inward normal of $e_b$ is positive on $a$, the parallelogram
\[
 J=\conv(e_b)+[0,ra]
\]
lies in $\conv S$. Put $t=|e_b|-1$. The existence of an ambiguous base pattern and the stated edge
bound give $t\ge1$.

Each $b$-parallel lattice line through $J$, other than the line of $e_b$, contains at least $t$
consecutive sites of $S\setminus e_b$. Along its sufficiently far positive ray, the length-$t$ words
are restrictions of the ambiguous base patterns. There are at most $t$ such patterns by
\eqref{eq:ambiguity-budget}, so the one-sided Morse--Hedlund theorem gives an eventually periodic
tail with period at most $t$ \cite[Theorem~2.2(ii)]{CK}. Taking a common multiple over the finitely many lines through $J$, and
discarding their initial segments, gives a half-strip $H_0\subset\RR$ with period $h_b=lb$, $l\ge1$.

To start the row extensions, we need consecutive sites on the $a$-rows, not just on the
$b$-lines. Choose a primitive integer normal $\nu$ with $\nu(a)=0$ and $\nu(b)>0$, and let $g_0$
be the initial endpoint of $e_b$. Removing the line of $e_b$ from $J$ and extending in the positive
$b$-direction gives
\[
 J^+=\{g_0+\alpha a+\beta b:0<\alpha\le r,\ \beta\ge0\}.
\]
Every sufficiently high $\nu$-section of $J^+$ is a semi-open interval of length $r$ in primitive
$a$-units, so it contains $r$ consecutive lattice points. Each such point lies on a $b$-line whose
intersection with $J$ has primitive length $t\ge1$ and hence contains a lattice point. The
difference of the two points is an integer multiple of the primitive vector $b$. At sufficiently
high levels this multiple is positive and exceeds the corresponding eventual-periodicity threshold.
There are only finitely many $b$-lines through $J$, so a single threshold $j_0$ ensures that
$H_0\cap\{\nu=j\}$ contains $r$ consecutive $a$-sites whenever $j\ge j_0$. No assumption that
$|\det(a,b)|=1$ is needed.

The period $pa$ transports the same $h_b$-equalities to $H_0+Mpa$ for every integer $M\ge0$. These
translates supply $r$ consecutive equality sites arbitrarily far along each $a$-row above $j_0$. The
same sites give equality for every positive multiple of $h_b$.

Let $d$ be the $\nu$-width of $S$. Choose a finite lattice parallelogram $F\subset\RR$ containing
$p$ consecutive $a$-sites on each of $d+1$ successive levels above $j_0$. Complete $a$ to a
unimodular basis to construct this finite set, then place it sufficiently far inside $\RR$. This
auxiliary choice of coordinates does not change the observation rectangle. Every $F+qh_b$, $q\ge0$,
remains in $\RR$.

A repeated $F$-pattern at indices $q<q'$ gives agreement between $x$ and $T^{(q'-q)h_b}x$ on
$F+qh_b$. On each of its rows, both arrays are $pa$-periodic on the row half-line in $\RR$.
Agreement at $p$ consecutive sites therefore gives agreement on the whole row, producing a band of
$d+1$ rows.

Suppose first that the lower $a$-parallel edge has primitive length $r$. We extend such a band
downwards to level $j_0$. On the next row, choose an $r$-site equality seed far enough along a
translate $H_0+Mpa$. Place the lower edge of $S$ there, with all remaining sites in the known band.
Generation of the final endpoint gives the next equality on the row. Iterating along $a$ gives a
tail of equalities, and the common period $pa$ extends them to the whole row half-line. The same
argument applies to each successive lower row.

A different gap for each band would not suffice. Let $N$ be the number of possible $F$-patterns.
Every disjoint index block $[m(N+1),m(N+1)+N]$ contains a repeat with gap in $\{1,\ldots,N\}$, so
one gap $D$ occurs in infinitely many blocks. For this fixed $D$, the corresponding bands have
heights tending to infinity, since $\nu(h_b)>0$. Their downward sweeps cover every row above $j_0$
and all concern the same difference $\Del{Dh_b}x$. Hence $Dh_b$ is a period on
$\RR\cap\{\nu\ge j_0\}$.

If the upper $a$-parallel edge has the smaller length $r$, one repeated band suffices. Using that
edge, the same generated-endpoint argument extends agreement upwards through every row above the
band. Its fixed gap again gives a period $Dh_b$ on a smaller two-ray region. In both cases $pa$
remains a period, and the recession cone is unchanged.
\end{proof}

\section{Half-strips and anchored compactness}\label{sec:compactness}
A \emph{half-strip} in a primitive direction $a$, based on a finite polygon $F$, is
$(F+\R_{\ge0}a)\cap\Z^2$. If $F$ is $\ZZ$-enveloped and $a$ is a factor direction, this is also the
union of the nonnegative integer $a$-translates of $F\cap\Z^2$. Indeed, every parallel lattice fiber
of $F$ has length at least one primitive unit and contains a lattice site.

\begin{lemma}[A periodic half-strip supplies a wedge]\label{lem:initial-wedge}
Assume that $Q_0$ has at least two factor directions, and let $a$ be primitive in one of them.
Suppose $Q_0z=0$ and $z$ agrees with a globally $a$-direction-periodic array on a forward
$a$-half-strip based on a $\ZZ$-enveloped polygon. Then $z$ has a period $h$ in direction $a$ on a
translated region $\ZZ+\R_{\ge0}a+\R_{\ge0}b$, where $b$ is the next factor direction on one side of
$a$.
\end{lemma}
\begin{proof}
Choose a common multiple $h$ of the reference period and the parallel factor vector of $Q_0$, with
its sign pointing along the forward half-strip. The difference $E=\Del h z$ vanishes on that
half-strip. As in Lemma~\ref{lem:globalize}, cancelling the parallel factor gives $Q_1E=0$, where
$Q_1$ is the nonempty product of the remaining factors. We use this recurrence to extend the zero
set into an adjacent wedge.

Complete $a$ to a unimodular basis and use integral coordinates $(\tau,\nu)$ with $\tau(a)=1$ and
$\nu(a)=0$. Reverse each remaining factor vector $p_j$ so that $\nu(p_j)>0$. Let
\[
 d=\sum_j\nu(p_j),\qquad C=\max_j\frac{-\tau(p_j)}{\nu(p_j)}.
\]
After multiplication by a Laurent unit, $Q_1$ has nonzero constant coefficient. Every other exponent
$u$ satisfies
\[
 1\le\nu(u)\le d,\qquad \tau(u)+C\nu(u)\ge0.
\]
These inequalities hold for every subset sum of the factor vectors, and thus for the actual support
even when coefficients cancel.

For some integers $M,s$, the half-strip contains $\{\tau\ge M,\ s\le\nu\le s+d-1\}$. To see this,
take a translate of $\ZZ$ in its base polygon. This zonotope has $\nu$-width $d$, and its factor
segment parallel to $a$ gives a lattice site on each integer level in the indicated range. Forward
$a$-translation then fills a tail on each of these finitely many levels, allowing a common threshold
$M$. Put
\[
 A=M+\max_{s\le j\le s+d-1}Cj.
\]
The recurrence of $Q_1$ propagates $E=0$ downwards through all integer $\nu$-levels in
\[
 W=\{\nu\le s+d-1,\ \tau+C\nu\ge A\}.
\]
The initial $d$ levels of $W$ lie in $\{\tau\ge M\}$ by the choice of $A$. At any lower level, every
nonconstant recurrence input has a larger $\nu$-value, at most $s+d-1$, and no smaller value of
$\tau+C\nu$. All these inputs are already zero, so the nonzero constant coefficient determines the
next level. Induction gives $E=0$ throughout $W$.

The recession rays of $W$ are $a$ and a primitive vector $b$ antiparallel to a factor vector
attaining the maximum in the definition of $C$. They are nonparallel. By this extremal choice of
slope, no factor direction lies strictly between them. An integer translate of $\ZZ$ sufficiently
far inside $W$ remains in $W$ after either recession ray is added. In particular, $\Del h
z=0$ on the resulting region, as required.
\end{proof}

Maximality gives a disagreement somewhere outside each agreement polygon. That site may move with
the polygon, and then compactness need not preserve the disagreement. We anchor a corner before
taking a limit and use the facet recurrence to locate a witness in a fixed finite set near that
corner.

\Needspace{14\baselineskip}
\begin{lemma}[Anchored agreement with a retained boundary failure]\label{lem:anchored}
Let $a$ be primitive in a factor direction of $Q_0$, let $H$ be a lattice half-plane whose boundary
is oriented by $a$ with $H$ on its left, and let $w\in X_c$ have a period in direction $a$. Suppose
that no shift of $c$ agrees with $w$ on any forward $a$-half-strip based on a $\ZZ$-enveloped
polygon. Then there are a fixed phase $\widehat w$ of $w$, a point $\theta\in X_c$, and a two-ray
region $\RR\subset H$ with Minkowski summand $\ZZ$ such that:
\begin{enumerate}[label=(\roman*)]
\item its recession rays are $-a$ and a nonparallel primitive factor direction $b$;
\item $\theta=\widehat w$ on $\RR$, but not on its next outward lattice layer $\RR^{(1)}$ along the $b$-edge;
\item $\RR^{(1)}\subset H$, and every finite subset of $H$ can be moved into $\RR$ by a sufficiently large negative multiple of $a$.
\end{enumerate}
\end{lemma}
\begin{proof}
\emph{Maximal agreement polygons.} After completing $a$ to an integral basis, call the forward
$a$-direction rightwards; we use these auxiliary coordinates only to describe the polygons.
Fix a lattice point $g_1$ on the boundary of $H$. Inductively choose a $\ZZ$-enveloped lattice polygon $F_i$, supported on this boundary, which contains $M_{i-1}$ and $[-i,i]^2\cap H$. For the first step, omit $M_{i-1}$. Such polygons are obtained by sufficiently large integral dilations of $\ZZ$, translated along their nondegenerate $a$-edge. Since $w\in X_c$, choose $u_i$ so that $T^{u_i}c=w$ on $F_i$. Among the $\ZZ$-enveloped polygons between $F_i$ and its forward half-strip on which these arrays agree, choose an inclusion-maximal polygon $M_i$.

To justify this choice, the half-strip fixes the transverse width. Along the finitely many
nonhorizontal edge directions, the longitudinal change is bounded by that width times a constant
depending only on the slopes. Thus if any right endpoint were unbounded, all right endpoints would
tend to infinity. The left boundary is fixed by $F_i$, while some lattice site of the half-strip is
a disagreement by hypothesis. This bounds all right endpoints, so there are only finitely many admissible
lattice polygons. In particular $M_i$ exists, and
\[
 M_{i-1}\subset F_i\subset M_i.
\]

\smallskip\noindent
\emph{Anchoring and convergence.} Translate the final endpoint of the oriented $a$-edge of $M_i$ to $g_1$, writing the translation as $-k_i a$. Set
\[
 P_i=M_i-k_i a,\qquad \theta_i=T^{u_i+k_i a}c.
\]
Then $\theta_i=T^{k_i a}w$ on $P_i$. The reference has a period parallel to $a$, so after passing to
a subsequence $T^{k_i a}w$ is one fixed array $\widehat w$.

The different anchor translations need not preserve the nesting of the $M_i$. Pass to a further
subsequence on which every primitive edge length is nondecreasing. This is possible for a sequence
in $\mathbb N^{2q}$: in each coordinate, stabilize a bounded infinite subsequence if there is one,
and otherwise take a strictly increasing subsequence. Now the anchored $P_i$ are nested. The
differences of successive cyclic edge-vector lists have nonnegative lengths and sum to zero, so
they bound a convex polygon, possibly a segment or point, anchored at the origin. Minkowski addition
adds both these lists and the anchored vertices. Each later $P_i$ is therefore the earlier polygon
plus this residual polygon, which contains the origin.

The $a$-edge lengths and the transverse widths tend to infinity. In the cyclic list following the
$a$-edge, let $b$ be the first direction with unbounded length before the opposite $a$-edge. Since the widths diverge, such a
direction exists; all intervening edge lengths stabilize. The increasing
union of the $P_i$ is a rational two-ray polyhedron $R_\infty$ with recession rays $-a,b$ and a
fixed finite boundary chain in $H$.

More explicitly, the supporting lines of that finite chain stabilize. Every other supporting normal
is negative on at least one of $-a,b$, so its support level tends to minus infinity. Precisely the
finite-chain inequalities remain, and every point satisfying them eventually lies in $P_i$. Thus the
union is the resulting closed polyhedron. Subtracting $\ZZ$ from the edge lists, or taking the
increasing union of the corresponding summands, also shows that $R_\infty$ retains a Minkowski
summand $\ZZ$.

Passing to a subsequence, we may assume that $\theta_i\to\theta$. Every fixed finite part of $R_\infty$
eventually lies in $P_i$, so $\theta=\widehat w$ there. Moreover $\theta\in X_c$, and hence
$Q_0\theta=0$ and every finite pattern of $\theta$ occurs in $c$.

\smallskip\noindent
\emph{A fixed boundary witness.} Write $\nu_b=\beta$ for the infinite $b$-facet of $R_\infty$, with primitive integer inward normal. For large $i$, the $b$-facet of $P_i$ lies on this same line. Let $P_i^{(j)}$ be obtained by replacing its inequality with $\nu_b\ge\beta-j$ and retaining all other inequalities. For each fixed $j$, these polygons retain a $\ZZ$ summand for large $i$: the moving edge has length tending to infinity, its neighbors lengthen, and its own length changes by a bounded amount.

An expansion by one lattice level need not leave lattice vertices. Let $\mu_-,\mu_+$ be the
primitive inward normals of the neighboring facets, and choose $m\ge1$ divisible by the fixed
nonzero integers $|\det(\mu_-,\nu_b)|$ and $|\det(\mu_+,\nu_b)|$. To see why this choice works,
an endpoint $v$ incident to the neighbor with normal $\mu$ moves according to
\[
 \begin{pmatrix}\mu\\ \nu_b\end{pmatrix}
 (v^{(j)}-v)=\begin{pmatrix}0\\-j\end{pmatrix}.
\]
By Cramer's rule, the displacement is integral at $j=m$. All other vertices are unchanged, so
$P_i^{(m)}$ is a lattice polygon. Intermediate expansions may have rational vertices, as allowed in
Lemma~\ref{lem:facet}. For large $i$, the final polygon is $\ZZ$-enveloped and strictly larger than
$P_i$. It still lies in the allowed translated half-strip: the relaxed normal is negative on $a$, so
its inequality does not define the half-strip, while every defining inequality of the half-strip was
retained.

Let $t$ be the primitive length of the $b$-face of $\ZZ$. For each $1\le j\le m$, take the first $t$
lattice points, in direction $b$, on the new facet of $P_i^{(j)}$. Their union $E$ is independent of
all sufficiently large $i$. Indeed, the initial endpoint is determined by the stabilized preceding
supporting line and the fixed equation $\nu_b=\beta-j$; only the far endpoint escapes.

For every sufficiently large $i$, the arrays $\theta_i$ and $\widehat w$ must disagree somewhere on
$E$. Otherwise apply Lemma~\ref{lem:facet} successively to $\delta_i=\theta_i-\widehat w$, which
satisfies $Q_0\delta_i=0$. Agreement on $P_i$, together with the first $t$ zeros on each successive
facet, would give agreement on $P_i^{(m)}$, contradicting the maximality of $M_i$ after undoing the
anchor translation.

This is now a condition on one fixed finite pattern. Convergence on $E$ means eventual equality of
entire $E$-patterns, so
\[
 \theta|_E\ne\widehat w|_E.
\]
Since $E\subset R_\infty^{(m)}$, one of the first $m$ outward layers contains a disagreement. Take
the first such layer and let $\RR$ be the preceding region. Then $\theta=\widehat w$ on $\RR$ but
not on $\RR^{(1)}$. The defining inequality of $H$ was never moved, so $\RR^{(1)}\subset H$. Every
other surviving inward normal is strictly positive on $-a$; hence a sufficiently large negative
$a$-translation moves any fixed finite subset of $H$ into $\RR$.
\end{proof}

\begin{lemma}[The selected point cannot acquire a period]\label{lem:nonperiodic}
In Lemma~\ref{lem:anchored}, suppose in addition that $w|_H$ is not fully periodic. Then $\theta$ is
nonperiodic.
\end{lemma}
\begin{proof}
We test a hypothetical period against the fixed reference, rather than the approximating arrays.
The phase shift is parallel to the boundary of $H$, so $\widehat w|_H$ is not fully periodic. Let
$a_*=la$, $l\ge1$, be a period of $\widehat w$, and suppose that $p\ne0$ is a period of $\theta$.
Given $g,g+p\in H$, subtracting $Na_*$ puts both points in $\RR$ for large $N$, whence
\[
 \widehat w(g)=\theta(g-Na_*)
 =\theta(g+p-Na_*)=\widehat w(g+p).
\]
This makes $p$ a regional period of $\widehat w|_H$. Together with $a_*$ it would give full regional
periodicity unless $p$ were parallel to $a$.

In the remaining case, primitivity gives $p=ra$ for some nonzero integer $r$. Both arrays then have
period $\lcm(l,|r|)a$. Transporting any point of $H$ into $\RR$ by a sufficiently large negative
multiple of this common period extends their agreement to all of $H$, contradicting the disagreement
on $\RR^{(1)}\subset H$.
\end{proof}

\section{Proof of the selection proposition}\label{sec:selection-proof}
\begin{proof}[Proof of Proposition~\ref{prop:selection}]
Fix $S$, $Q_0$, and the periodic reference $w$ from Lemmas~\ref{lem:generating}
and~\ref{lem:reference}. The product $Q_0$ has at least two factors, since a single factor would
give a period of $c$. Its directions include $a$, and $S$ has two opposite $a$-parallel edges and
satisfies \eqref{eq:edge-budget}. We keep this reference fixed throughout. We divide into two cases, depending on
whether one of its forward half-strips already occurs in a shift of $c$.

\emph{A matching half-strip occurs.} Suppose that some shift $z=T^uc$ agrees with $w$ on a forward $a$-half-strip based on a $\ZZ$-enveloped polygon. By Lemma~\ref{lem:initial-wedge}, $z$ has a period $h$ parallel to $a$ on a translate $\RR_0$ of $\ZZ+\R_{\ge0}a+\R_{\ge0}b$, where $a,b$ are neighboring factor directions.

We can enlarge the region towards a half-plane by adding successive factor rays on the side of $b$
until reaching $-a$. Each new ray is added to the entire preceding region, with the earlier rays
retained. Adjacency in the cyclic list makes each step remove precisely the inequality of the
current infinite facet. The period $h$ cannot persist to the final half-plane, whose boundary is
parallel to $a$: Lemma~\ref{lem:globalize} would give a global period of $z$ and hence of $c$.

At the first enlargement on which the period fails, we relax the moving inequality through successive integer levels. These finite relaxations exhaust the enlargement, so there is a last region on which
$\Del h z=0$ and a next layer on which it is nonzero. The points $z,T^hz\in X_c$ agree on the former
region and not on the latter layer. Lemma~\ref{lem:boundary-ambiguity} gives ambiguity in that
boundary direction, say $b'$. Its inward normal is positive on $a$, so Lemma~\ref{lem:regional}
supplies periods parallel to $a,b'$ on a smaller two-ray region. We take $\theta=z$, which is
nonperiodic as a shift of $c$.

\emph{No matching half-strip occurs.} We apply Lemma~\ref{lem:anchored} to $w$ and $H$. It gives a point $\theta\in X_c$ agreeing with a fixed phase $\widehat w$ on a two-ray region with recession rays $-a,b$, and disagreeing on the next outward layer. By Lemma~\ref{lem:nonperiodic}, this point is nonperiodic. Lemma~\ref{lem:boundary-ambiguity} gives ambiguity along $b$. Agreement with the reference supplies a regional period pointing along $-a$, and the inward normal of the $b$-edge is positive on $-a$. Lemma~\ref{lem:regional} therefore gives two independent periods on a smaller region.

In either case, choose the signs of the two periods along the recession rays.
Formula~\eqref{eq:background} extends the restriction to a single fully periodic background $B$. The
region contains a translate $K$ of its rational recession cone, of positive aperture less than
$\pi$, and $\theta=B$ on $K\cap\Z^2$. Finally, $\theta\in X_c$ gives $P_\theta(k,n)\le P_c(k,n)$ on
the original rectangle.
\end{proof}

\appendix
\section{Relation to previous work}\label{app:sources}
We use two structural results from the literature: Kari--Szabados, Corollary~13 \cite{KS}, and Colle,
Theorem~1.9 \cite{Colle}, both quoted without proof. The one-dimensional arguments use the bi-infinite
and one-sided forms of the Morse--Hedlund theorem recorded by Cyr--Kra, Theorem~2.2(ii)--(iii)
\cite{CK}; see also the original paper \cite{MH}.

The proof of Lemma~\ref{lem:affine-ann} follows the affine argument in \cite[Lemma~5]{KS}, while
Lemma~\ref{lem:coprime} uses the elimination argument of \cite[Lemma~8]{KMalg}, with the finite-image
hypothesis stated explicitly. For integral configurations, the directional structure underlying
Lemma~\ref{lem:gcd} is contained in \cite[Corollaries~28--29]{KS}. Our proof starts from the product
annihilator and also covers the finite-valued complex configurations in Theorem~\ref{thm:wedge}. The
affine pattern count is related to \cite[Lemmas~33 and~36]{KS}. Products of directional layers also
occur in the four-dot argument of Kari--Moutot \cite[proof of Theorem~11]{KMalg}. The contribution of
Sections~\ref{sec:observable}--\ref{sec:count} is the construction of a quadratic observable with a
small stencil and the proof that its fitting translates are independent modulo the affine functions on
the same rectangle.

We use Colle's half-plane formulation of one-sided nonexpansiveness \cite[Section~1]{Colle}; Cyr--Kra
\cite[Definition~3.1 and Lemma~3.2]{CK} give the finite-extension formulation and its relation to
half-plane ambiguity. Colle \cite[Proposition~2.14]{Colle}, following Kari--Szabados, obtains the
tangent-period conclusion from a periodic decomposition; Lemma~\ref{lem:globalize} proves it directly
from the product annihilator. The period-difference construction in Lemma~\ref{lem:initial-wedge} also
appears in Colle \cite[proof of Lemma~4.6, Claim~4.7]{Colle}; our proof writes out the propagation
into the wedge in integral coordinates.

Lemma~\ref{lem:reference} is based on the ambiguity argument of Cyr--Kra \cite[Lemma~2.24]{CK} and the
outward-extension argument from the proof of their Proposition~4.8, as presented by Colle
\cite[Proposition~2.13]{Colle}. Since the full Cyr--Kra proposition assumes $P_c(k,n)\le kn/2$, we
prove the extension step directly from the edge ambiguity bound at the $kn$ threshold. Related
bounded-period propagation appears in Colle--Garibaldi \cite[Lemmas~4.11--4.13]{CG}.

The argument in Lemma~\ref{lem:regional} follows Colle's Lemma~4.2 \cite{Colle}, with the hypotheses
arranged in the form needed here. In Proposition~\ref{prop:selection}, we adapt the two-case
region-selection argument in his proof of Lemma~4.6 to obtain a periodic region in a configuration
that remains nonperiodic. For this purpose, the later conclusions about the angular positions of the
nonexpansive directions, as well as Theorem~1.10 and its constant decomposition-order hypothesis, are
not needed.

Related periodic-region arguments appear in Colle--Garibaldi \cite[Sections~5--6]{CG} under an
alphabet-sensitive complexity hypothesis. The more recent modular structure theorem
\cite[Theorem~1.2]{CG26} concerns reductions of periodic components on pairs of half-planes, and the
proof of Theorem~3.5 in that paper again uses the two-case construction. Here we keep a single orbit
point over the original alphabet and prove its nonperiodicity directly, without imposing the earlier
alphabet-sensitive hypothesis. The short-side-four result provides further context, though it is not
actually used in the proof.

The anchored construction in Lemma~\ref{lem:anchored} is based on Colle \cite[proof of
Lemma~3.5]{Colle}, with an extra boundary argument to retain a disagreement in the limit. We shift the
approximants by $u_i+k_i a$, reducing only the periodic reference phase to a fixed residue. After
passing to a componentwise monotone subsequence of edge lengths, the anchored polygons are nested by
the residual Minkowski-sum argument. Outward layers are formed by relaxing signed supporting
inequalities, so that each new region contains the preceding one; an unsigned distance restriction on
the entire region would clip it. Lemma~\ref{lem:facet} propagates agreement from the fixed finite set
$E$, supplying the boundary control that keeps disagreement from escaping under compactness. Finally,
Lemma~\ref{lem:nonperiodic} avoids the two ray-limit arguments by transporting a hypothetical period
of $\theta$ to the reference half-plane and using the retained disagreement, without transferring a
period backwards from a limit.

\section{Boundary cases and examples}\label{app:examples}
Finite image is needed in Lemma~\ref{lem:coprime}, although boundedness suffices for the later
arithmetic-progression arguments applied to $c$ and $A_i$. For example, if
$1,\alpha/(2\pi),\beta/(2\pi)$ are rationally independent, then $e^{i(\alpha x+\beta y)}$ is bounded
and has two coprime univariate annihilators, but no period. Its image is infinite.

The single defect has exactly $kn+1$ patterns but no Laurent annihilator by Lemma~\ref{lem:cone}:
it gives sharpness without satisfying the algebraic hypothesis of Theorem~\ref{thm:wedge}. A lifted
aperiodic sequence has one period, though not two independent periods, and is excluded by the
nonperiodicity hypothesis.

An example meeting the algebraic and wedge hypotheses is
\[
 c(x,y)=\mathbf1_{x\ge0}+\mathbf1_{y\ge0}-\mathbf1_{x+y\ge0}\in\{0,1\}.
\]
It is annihilated by $\Del{(1,0)}\Del{(0,1)}\Del{(1,-1)}$ and is $1$ on the first quadrant. It is
nonperiodic: for $y>0$, the zero interval in its row is $[-y,-1]$. Comparing these intervals at
sufficiently large heights under a hypothetical period $(a,b)$ forces $a=0$ from the right endpoints
and then $b=0$ from the left endpoints. Thus the wedge hypotheses allow nonperiodic binary
configurations with three factor directions, including a negatively sloped direction.

Section~\ref{sec:count} covers every positive rectangle size, including $L_R=0$, in the original
coordinates; the auxiliary bases used in the selection do not alter the window. The cone argument does not extend to dimension three, since two half-spaces can intersect in a set containing a line.

\begingroup
\raggedright
\begin{thebibliography}{10}
\bibitem{Colle}
C. F. Colle,
\emph{On periodic decompositions, one-sided nonexpansive directions and Nivat's conjecture},
Discrete and Continuous Dynamical Systems \textbf{43} (2023), no.~12, 4299--4327.
\href{https://doi.org/10.3934/dcds.2023088}{doi:10.3934/dcds.2023088}.

\bibitem{CG}
C. F. Colle and E. Garibaldi,
\emph{An alphabetical approach to Nivat's conjecture},
Nonlinearity \textbf{33} (2020), no.~7, 3620--3652.
\href{https://doi.org/10.1088/1361-6544/ab7638}{doi:10.1088/1361-6544/ab7638}.

\bibitem{CG26}
C. F. Colle and E. Garibaldi,
\emph{A modular structure theorem for minimal periodic decompositions and periodicity of
configurations with $P_\eta(4,n)\le4n$},
preprint, \href{https://arxiv.org/abs/2606.10193}{arXiv:2606.10193} (2026).

\bibitem{CK}
V. Cyr and B. Kra,
\emph{Nonexpansive $\Z^2$-subdynamics and Nivat's conjecture},
Transactions of the American Mathematical Society \textbf{367} (2015), no.~9, 6487--6537.
\href{https://doi.org/10.1090/S0002-9947-2015-06391-0}{doi:10.1090/S0002-9947-2015-06391-0}.

\bibitem{CKshort}
V. Cyr and B. Kra,
\emph{Complexity of short rectangles and periodicity},
European Journal of Combinatorics \textbf{52} (2016), 146--173.
\href{https://doi.org/10.1016/j.ejc.2015.10.003}{doi:10.1016/j.ejc.2015.10.003}.

\bibitem{KMalg}
J. Kari and E. Moutot,
\emph{Nivat's conjecture and pattern complexity in algebraic subshifts},
Theoretical Computer Science \textbf{777} (2019), 379--386.
\href{https://doi.org/10.1016/j.tcs.2018.12.029}{doi:10.1016/j.tcs.2018.12.029}.

\bibitem{KM}
J. Kari and E. Moutot,
\emph{Decidability and periodicity of low complexity tilings},
Theory of Computing Systems \textbf{67} (2023), 125--148.
\href{https://doi.org/10.1007/s00224-021-10063-8}{doi:10.1007/s00224-021-10063-8}.

\bibitem{KS}
J. Kari and M. Szabados,
\emph{An algebraic geometric approach to Nivat's conjecture},
Information and Computation \textbf{271} (2020), article~104481.
\href{https://doi.org/10.1016/j.ic.2019.104481}{doi:10.1016/j.ic.2019.104481}.

\bibitem{MH}
M. Morse and G. A. Hedlund,
\emph{Symbolic dynamics},
American Journal of Mathematics \textbf{60} (1938), no.~4, 815--866.
\href{https://doi.org/10.2307/2371264}{doi:10.2307/2371264}.

\bibitem{Sz}
M. Szabados,
\emph{Nivat's conjecture holds for sums of two periodic configurations},
in \emph{SOFSEM 2018: Theory and Practice of Computer Science},
Lecture Notes in Computer Science, vol.~10706, Springer, 2018, pp.~539--551.
\href{https://doi.org/10.1007/978-3-319-73117-9_38}{doi:10.1007/978-3-319-73117-9\_38}.
\end{thebibliography}
\endgroup
\end{document}
